%% !TEX root = ../manual_gaolv5.tex
% Chapter 11.
%% Chapter 11: Accuracy of the operations. Sources: gaol.tex of f5bf180 5881-6006
%% (chapter "Accuracy of the operations"); doc/accuracy.md (whole).
\chapter{Accuracy of the operations}\label{chap:accuracy}
%%=================================
\cindex{accuracy!of the operations}\cindex{tightness}%

The containment property (\pxref{sec:containment}) says that the result of
an operation encloses its exact value; it says nothing of how closely.
IEEE 1788-2015 \cite{IEEE1788:2015} requires an implementation to document
the \dfn{tightness} of each of its interval operations, dividing the
possible inputs into ranges and stating the tightness achieved in each
(12.10.3). This chapter does so for \code{interval}, the inf-sup type of
GAOL on doubles, from the algorithm of each operation. The tests of
\file{tests} check each statement on random and special arguments
(\pxref{sec:test-programs}), and we give, with each operation, the largest
distance they found, or the distance they accept where the text says no more.\newinvfive

\medskip
IEEE 1788-2015 defines three accuracy modes (12.10.1), which
Chapter~\ref{chap:concepts} presents (\pxref{sec:accuracy-modes}). For an
operation $f$, its exact value $f^*$ and an input box $x$:
\begin{itemize}
\item the result is \emph{tightest} when $f(x)$ is the hull, in doubles, of
  $f^*(x)$: each bound is the nearest double on its side of the exact bound;
\item it is \emph{accurate} when $f(x)$ encloses $f^*(x)$ and is within
  \[\text{nextOut}(f_{\text{tightest}}(\text{nextOut}(x))),\]
  nextOut moving each bound one double outward: roughly, $f(x)$ is no wider
  than the tightest result for an input one double wider, plus one double on
  each side;
\item it is \emph{valid} when $f(x)$ encloses $f^*(x)$.
\end{itemize}

\noindent
The standard requires the basic operations (\ieee{neg}, \ieee{add}, \ieee{sub},
\ieee{mul}, \ieee{div}, \ieee{recip}, \ieee{sqr}, \ieee{sqrt}, \ieee{fma}),
the integer functions and \code{abs}, \code{min} and \code{max} to be
tightest; it requires the other elementary functions, the reverse functions
and the recommended functions to be valid, and recommends them to be
accurate (12.10.2). Below, ``within $k$ doubles'' means that each bound is
at most $k$ doubles beyond the tightest one.

\medskip
GAOL bounds its elementary functions with CORE-MATH \cite{CoreMath:2022},
which is correctly rounded in the rounding direction in effect
(\pxref{sec:math-libraries}).\cindex{CORE-MATH!accuracy}
Computing in the upward rounding it keeps, GAOL takes the value at a bound
as the upper bound, with nothing to add, and the double below it as the
lower one. The upper bound is the tightest one, and so is the lower bound unless the exact
value is a double, which the operations give exactly where it happens:
$\log(1)=0$ and $\sin(0)=0$ are exact, and $\arcsin(1)$ is the tightest
enclosure of $\pi/2$. Where an algorithm evaluates the function at the
bounds of the interval, the bounds of the result are therefore the tightest
ones.

\medskip
The first section of this chapter states the conditions under which these
statements hold. The second gives the tightness of each operation in a
table, then, family by family, the algorithm of each operation and the
largest distance from the tightest bounds that the tests found. The third
section explains why the bounds are the same on every machine, and where
the sign of a zero bound differs between the kernels.

\section{The conditions}\label{sec:accuracy-conditions}
%%----------------------------------------------------
\cindex{accuracy!conditions}%

The statements of this chapter hold for GAOL as the three builds of
Chapter~\ref{chap:installation} make it: with the CORE-MATH of
\file{3rd/math-core}, compiled as the builds compile it, and with the code
that includes the headers of GAOL compiled with the flags of interval
arithmetic (\pxref{sec:compiling}). They hold for the SSE2 intervals, the
NEON intervals\newinvfive{} and the FPU intervals alike
(\pxref{sec:kernels}), and for the AVX-512 path (\pxref{sec:avx512}), whose
bounds are those of the SSE2 path, but for the sign of a zero bound of
\code{x * d} (\pxref{sec:accuracy-same-bounds})

\medskip
They hold whatever the rounding direction the calling code left
(\file{tests/\allowbreak rounding\_direction.cpp}), and with the floating-point
exceptions masked, as a program starts: an operation that raises an enabled
exception stops the program rather than give its bounds
(\pxref{sec:fp-exceptions}). They hold on every architecture and with every
compiler alike, CORE-MATH giving the same bits everywhere: there is no other
mathematical library to build GAOL with (\pxref{sec:math-libraries}, and
Chapter~\ref{chap:changes} for what differs from GAOL~4). The intervals of
floats, \code{gaol::intervalf} and \code{gaol::interval2f}, unfinished,
which only a developer of GAOL compiles (\file{gaol/gaol\_config.h}), are
not covered.

\section{The tightness of each operation}\label{sec:accuracy-table}
%%-----------------------------------------------------------------
\cindex{accuracy!table of the operations}%

Table~\ref{tab:accuracy-tightness} gives the tightness of each operation of
\code{interval}, in the families of IEEE 1788-2015, from its algorithm; the
subsections that follow give the algorithm of each operation and what the
tests found. The tables numbered 9.1, 10.1, 10.3, 10.4, 10.5 and 10.7 that this
chapter names are those of IEEE 1788-2015.

\begingroup\renewcommand{\arraystretch}{1.1}\small
\begin{longtable}{L{4.8cm}L{9.2cm}}
  \caption{The tightness of each operation of \code{interval}}%
  \label{tab:accuracy-tightness}\\
  \hline
  \multicolumn{1}{c}{Operations} & \multicolumn{1}{c}{Tightness}\\
  \hline\hline
  \endfirsthead
  \multicolumn{2}{l}{\textit{Table~\thetable, continued}}\\*
  \hline
  \multicolumn{1}{c}{Operations} & \multicolumn{1}{c}{Tightness}\\
  \hline\hline
  \endhead
  \hline
  \endfoot
  \multicolumn{2}{l}{\textit{Basic operations (Table~9.1): tightest
    required}}\\*
  \code{-x}, \code{+}, \code{-}, \code{*}, \code{/}, \code{inverse},
  \code{sqr} & tightest\\
  \code{sqrt} & tightest, whatever the rounding of the square root of the C
  library\\
  \code{fma} & tightest\\
  \hline
  \multicolumn{2}{l}{\textit{Integer and absmax functions: tightest
    required}}\\*
  \code{floor}, \code{ceil}, \code{integer}, \code{sign}, \code{trunc},
  \code{round\_ties\_to\_even}, \code{round\_ties\_to\_away}, \code{abs},
  \code{min}, \code{max} & tightest, and exact, in every rounding
  direction\\
  \hline
  \multicolumn{2}{l}{\textit{Powers, exponentials and logarithms: valid
    required, accurate recommended}}\\*
  \code{pow(x, n)}, integer \var{n} (\ieee{pown}) & the tightest bounds, for a base of either sign, or one double beyond where the power is within $2K$ of a
  double, relatively, $K$ being the bound of the relative error of the power
  ($K=2^{-98}$ for $n=3$, about $2^{-62}$ at most, at $n=-2^{31}$; no such case for $n=2$, which is
  \code{sqr(x)}); this holds below $2^{-968}$ and beyond the largest double
  too; the bounds are exact where the power is a double; the same bounds
  with the SSE2, the NEON and the FPU intervals. Beyond the \code{int}s:
  tightest where $n$ is a double, within three doubles otherwise\\
  \code{pow(x, y)} (\ieee{pow}) & tightest, at the corners of the box and
  at their limits where the base is $0$ or $+\infty$ or the exponent
  infinite; for a degenerate integer exponent within the \code{int}s, the
  integer power of the row above, one double beyond the tightest where it
  is within twice the error of its products
  (\code{pow([0x1.fffffffffffffp+1], [7])} has its lower bound one double
  below the tightest)\\
  \code{exp}, \code{exp2}, \code{exp10}, \code{log}, \code{log2},
  \code{log10} & tightest; exact where the value is a double ($\exp(0)$,
  $\log(1)$, $2^n$, $10^n$ for $0\leq n\leq22$\ldots)\\
  \hline
  \multicolumn{2}{l}{\textit{Trigonometric and hyperbolic functions: valid
    required, accurate recommended}}\\*
  \code{cos}, \code{sin} & tightest where the function is evaluated at the
  bounds, at every magnitude; $-1$ and $1$ exactly at the extrema;
  \itv{-1}{1} from the width $2\underline{\pi}$ on, $\underline{\pi}$ being
  the double below $\pi$\\
  \code{tan} & tightest between two poles, at every magnitude;
  \itv{-\infty}{+\infty} exactly when a pole is within $x$, but for the
  widths above $\underline{\pi}$ and below $\pi$\\
  \code{asin}, \code{acos}, \code{atan}, \code{atan2} & tightest; exactly
  0, or the tightest enclosures of $\pm\pi/4$, $\pm\pi/2$ and $\pi$, where
  the value is one of them\\
  \code{sinh}, \code{cosh}, \code{tanh}, \code{asinh}, \code{acosh},
  \code{atanh} & tightest\\
  \hline
  \multicolumn{2}{l}{\textit{Recommended functions (Table~10.5)}}\\*
  \code{nth\_root(x, n)}, $n>0$ & $n=1$, 2, 3: tightest. Otherwise
  the tightest bounds or up to 2 doubles beyond, proved with
  integer powers; beyond the \code{unsigned int}s, CORE-MATH's \code{pow}
  at the corners, within one double\\
  \code{nth\_root(x, q)}, $q<0$ & valid: the root above, then the
  tightest division, up to 4 doubles beyond the tightest bounds\\
  \code{expm1}, \code{exp2m1}, \code{exp10m1}, \code{log1p},
  \code{log2p1}, \code{log10p1}, \code{hypot}, \code{rsqrt} & tightest\\
  \code{compoundm1} & tightest but in a case out of $10^{17}$ or so, where
  a bound is a double beyond\newinvfive\\
  \code{sinpi}, \code{cospi}, \code{tanpi}, \code{asinpi}, \code{acospi},
  \code{atanpi}, \code{atan2pi} & tightest, at every magnitude\\
  \hline
  \multicolumn{2}{l}{\textit{The error functions, which IEEE 1788-2015 does
    not name}}\\*
  \code{erf}, \code{erfc} & tightest; exact at 0 and at the
  infinities\newinvfive\\
  \hline
  \multicolumn{2}{l}{\textit{Reverse functions (Table~10.1): valid
    required, accurate recommended}}\\*
  \code{sqrt\_rel}, \code{nth\_root\_rel}, \code{acosh\_rel},
  \code{asinh\_rel}, \code{atanh\_rel} & accurate\\
  \code{nth\_root\_rel(c, p, x)}, $p<0$ & valid: the root of each part
  of $c$ and then its inverse, up to four doubles beyond the
  tightest\newinvfive\\
  \code{invabs\_rel} & tightest\\
  \code{max\_rel}, \code{min\_rel}, \code{sign\_rel}, \code{floor\_rel}
  (not in Table~10.1) & tightest\\
  \code{asin\_rel}, \code{acos\_rel}, \code{atan\_rel} & accurate, or one
  double beyond; a bound of $x$ beyond $2^{52}$ is kept\\
  \code{div\_rel}, \code{\%} & accurate; tightest where the bounds of the
  quotients are doubles\\
  \code{pow\_rel}, \code{pow\_exponent\_rel} & tightest, but for a bound of
  $x$ next to the solutions whose power is within one double of $c$, which
  is kept: the tightest result for $x$ one double wider\newinvfive\\
  \code{atan2\_rel\_y}, \code{atan2\_rel\_x} & accurate; within two doubles
  of the tightest bounds where these are not empty\\
  \code{cancel\_minus}, \code{cancel\_plus} & tightest\\
  \hline
  \multicolumn{2}{l}{\textit{Set operations, constructors, numeric and
    boolean functions}}\\*
  \code{\&}, \code{|}, \code{interval(l, u)} & tightest (exact)\\
  \code{textToInterval(s)}, \code{operator>{}>} & tightest for the literals
  of IEEE 1788-2015; valid for expressions, computed with the operations
  above\\
  \code{midpoint()}, \code{width()}, \code{rad()} & as 12.12.8 requires\\
  \code{mag()}, \code{mig()}, \code{smig()} & exact\\
  \code{mid()}, \code{hausdorff()} & tightest\\
  relations, \code{overlap()} & exact, as Tables~10.3, 10.4 and 10.7\\
  \code{operator<{}<} & \code{hexa}: exact, read back as the same interval.
  Decimal formats: valid, and the tightest with the digits asked for where
  the C library rounds to nearest as it should. \element{width} and
  \element{center}: $\itv{c-w}{c+w}$ contains $x$, $w$ being the smallest radius
  that does so. \code{exact\_string()}: exact, read back bit for bit\\
\end{longtable}
\endgroup

\begin{Rem}\label{rem:accuracy-before}
Before GAOL v5, the elementary functions but the hyperbolic ones were
correctly rounded to nearest and moved one double outward, which is at most one double beyond the
tightest bound: accurate, and not tightest. The hyperbolic functions were
taken from the mathematical library of the system, whose values are
sometimes further than one double from the exact ones, and millions of
doubles away for the \code{acosh} of some of them.\newinvfive
\end{Rem}

\subsection{Basic operations}\label{sec:accuracy-basic}

The basic operations of Table~9.1 of the standard must be tightest. Each is
computed from the bounds, rounded upward once (\pxref{sec:kernels}), and
each is tightest; the tests find them so.

\begin{itemize}
\item[$\bullet$] \textbf{Negation.} \code{-x} exchanges the bounds: exact.
\item[$\bullet$] \textbf{Addition and subtraction.} \code{x + y},
  \code{x - y}, \code{+=}, \code{-=}, and their forms with a double, are the
  sums of the bounds rounded upward, the lower bound being stored negated:
  tightest.
\item[$\bullet$] \textbf{Multiplication.} \code{x * y}, \code{*=}, and the
  forms with a double, are the products of the bounds chosen by their signs,
  rounded upward; a zero bound times an infinite one counts as $0$:
  tightest.
\item[$\bullet$] \textbf{Division.} \code{x / y}, \code{/=}, and the forms
  with a double, are the quotients of the bounds chosen by their signs,
  rounded upward; the hull of the quotients when $y$ contains $0$, and the
  empty set when $y$ is $\itv{0}{0}$: tightest.
\item[$\bullet$] \textbf{Inverse.} \code{inverse(x)} and \code{x.inverse()}
  divide $1$ by the bounds, rounded upward: tightest.
\item[$\bullet$] \textbf{Square.} \code{sqr(x)} squares the bounds, rounded
  upward, with $0$ as lower bound when $x$ contains $0$: tightest.
\item[$\bullet$] \textbf{Square root.} \code{sqrt(x)} works on
  $x\cap\itv{0}{+\infty}$. The upper bound $u$ is the square root rounded
  upward, from the square root of the C library checked with a product
  rounded the other way and moved one double when needed; the lower bound
  is $u$ when $u\cdot u$ is exactly the bound of $x$, and the double below
  $u$ otherwise: tightest, whatever the rounding of the square root of the
  C library.
\item[$\bullet$] \textbf{Fused multiply-add.} In \code{fma(x, y, z)}, the
  product $x\cdot y$ is bilinear, so it ranges over the hull of the
  products of the bounds, and $x\cdot y+z$ ranges from the least product
  plus $\leftBound{z}$ to the greatest plus $\rightBound{z}$. Each corner
  is rounded once by \code{std::fma}, upward for the upper bound and
  downward for the lower one, computed as $-\mathrm{fma}(-a, b, -c)$. A
  bound $0$ times an infinite one counts as $0$.\newinvfive{} Tightest,
  over 200\,000 boxes against exact rational arithmetic.
\end{itemize}

\begin{Rem}
The result of \code{std::fma} goes through \code{rnd\_keep()} before it is
negated. Without it, GCC folded $-\mathrm{fma}(-a, b, -c)$ into
$\mathrm{fma}(a, b, c)$, a single instruction rounded upward, and the lower
bound came above the exact value, one double above the tightest lower bound.
\end{Rem}

\subsection{Integer and absmax functions}\label{sec:accuracy-integer}

The integer functions, \code{abs}, \code{min} and \code{max} must be
tightest. None of them rounds: each returns a double that is an integer, or
a bound of its argument, in every rounding direction.

\begin{itemize}
\item[$\bullet$] \textbf{Floor and ceiling.} \code{floor(x)} and
  \code{ceil(x)} apply to each bound: exact.
\item[$\bullet$] \textbf{Integers of an interval.} \code{integer(x)}, which
  the standard does not name, is the hull of the integers in $x$,
  $\itv{\lceil\leftBound{x}\rceil}{\lfloor\rightBound{x}\rfloor}$: exact.
\item[$\bullet$] \textbf{Sign, truncation and roundings.} \code{sign(x)},
  \code{trunc(x)}, \code{round\_ties\_to\_even(x)} and
  \code{round\_ties\_to\_away(x)} take the value at each bound, each
  function being non-decreasing, as \code{ceil} and \code{floor} above, and
  $\emptyset$ for $\emptyset$, which the four test: GAOL holds the empty
  interval as two NaN bounds, which \code{trunc} and the two roundings send
  to themselves, where \code{sign} would give $\itv{0}{0}$, a NaN comparing
  false both to $0$ and above it. \code{trunc} and
  \code{round\_ties\_to\_away} are \code{std::trunc} and \code{std::round},
  which the C++ standard defines without reading the rounding direction,
  and \code{round\_ties\_to\_even} reads the bits
  (\file{gaol/math/gaol\_roundeven.h}), where \code{std::nearbyint} and
  \code{std::rint} would round in the direction in effect, upward in
  GAOL.\newinvfive{} Tightest and exact; the tests find the same result in
  the four rounding directions over 22\,407 doubles: the halfway values,
  the whole numbers and the doubles on either side of them.
\item[$\bullet$] \textbf{Absolute value.} \code{abs(x)} takes the magnitudes
  of the bounds, and $0$ as lower bound when $x$ contains $0$: exact.
\item[$\bullet$] \textbf{Minimum and maximum.} \code{min(x, y)} and
  \code{max(x, y)} apply to the bounds: exact.
\end{itemize}

\subsection{Powers, exponentials and logarithms}\label{sec:accuracy-powers}

These functions must be valid, and the standard recommends them to be
accurate. The exponentials and the logarithms are tightest, being
CORE-MATH's functions at the bounds; the integer power is tightest but one
double beyond in rare cases, and the real power is tightest at the corners
of its box.

\paragraph{The integer power.}\findex{pow}%
\code{pow(x, n)} for an integer $n$ of any type, and \code{pow(x, y)} and
\code{pow(x, d)} for a degenerate integer $y$ or an integer double $d$,
compute \ieee{pown} (\code{gaol\_pown()}, \code{gaol\_uipow()},
\code{gaol\_pown\_ll()}), as does \code{gaol\_ieee1788::pown(x, n)}
(\pxref{sec:powers}). Within the \code{int}s, $n=2$ is \code{sqr(x)}, the
squares of the bounds rounded outward.\newinvfive{} For any other $n$, the
powers of the two corners of $x$ where $x^n$ is the least and the greatest
are computed at once, in the two lanes of an SSE2 register with the FMA
instructions or of a NEON register, or of a pair of doubles otherwise (the
FPU intervals, and SSE2 without FMA), with the same operations in the same
order.\newinvfive{} The power is a binary exponentiation from exact
products: each product is rounded, and its rest, which \code{fma} gives
exactly, is carried along, so that the power is $h+l$ within $K\,|h|$ of
it, with $K=(8nL+16)\cdot2^{-104}$ for an $n$ of $L$ bits, a bound proved
beforehand (\file{gaol/core/gaol\_interval\_core.cpp}). Each bound is then
rounded once, $h+l+K|h|$ upward and $h+l-K|h|$ downward, and $K$ is left
out where every rest is $0$, $x^n$ being a double. The power $x^{-n}$ is
$1/(h+l)$: the inverse of $h$ rounded, corrected with the rest of
$1-q(h+l)$ to the second order and rounded once the same way; $x^{-1}$ is
the inverse. A corner whose powers leave $\itv{2^{-968}}{2^{960}}$ is
computed alone, scaled to $[1, 2)$, its power kept below $2^{501}$ by
powers of two, then scaled back with one rounding in the direction of the
bound, beyond the largest double and among the subnormals too. The bounds
$0$ and $\pm\infty$ give the limits of $x^n$:
$\itv{\mathrm{mag}(x)^n}{+\infty}$ where $x$ holds $0$ for an even $n<0$,
$\itv{-\infty}{+\infty}$ where it holds $0$ inside for an odd one, and
$\emptyset$ for $\itv{0}{0}$.

Beyond the \code{int}s, CORE-MATH's \code{pow} is called at the corners,
as for the real power below, the sign coming from the parity of $n$; for an
$n$ of more than 53 bits that is no double, $x^h\cdot x^l$ is rounded
outward, with $n=h+l$, $h$ a double and $|l|<2^{11}$; beyond 64 bits, $n$
is taken as $\pm(2^{63}-1)$ or $\pm(2^{63}-2)$, of its sign and its parity,
every power of a double but $0$, $\pm1$ and $\pm\infty$ being beyond the
doubles there. GAOL~4 gave $\itv{-\infty}{+\infty}$ beyond the
\code{unsigned int}s, and beyond the \code{int}s for \code{pow(x, y)};
within them, it rounded $1/x^n$ twice, and rounded the products outward, as
it computed every power, where a bound was $0$ or infinite or its power
beyond the doubles.\newinvfive

The result is tightest, or one double beyond where $x^n$ is within $2K$ of
a double, relatively: never for $n=2$, which is \code{sqr(x)}; $K$ is
$2^{-98}$ at $n=3$, $2^{-88}$ at $n=1000$, $2^{-65}$ at $n=2^{31}-1$ and
about $2^{-62}$ at $n=-2^{31}$, where one bound in about 400 is one double
beyond. It is exact where $x^n$ is a double. Beyond the \code{int}s it is
tightest where $n$ is a double, and within three doubles otherwise.

The tests find the tightest bounds at 5\,000 doubles drawn at random, with
exponents from $-30$ to $30$ and over all the doubles, subnormal results and
overflows included, for $n$ from 2 to 9, 16, 31 to 33 and $-1$ to $-4$,
$-7$, $-16$ and $-33$, and over 5\,000 intervals drawn so, for $n$ from 2
to 5, 12, 31 and $-1$ to $-3$, $-6$, $-17$, $x$ holding $0$ included,
against the exact powers; the tightest over 4\,000 intervals near the ends
of the doubles (a bound $0$ or infinite, powers below $2^{-1000}$ or above
$2^{1100}$); the exact value at the 8\,679 powers that are doubles; the
tightest at 13 powers for $n$ from $2^{24}+1$ to $2^{32}-1$, and at 6
corners computed alone, against mpmath. Beyond the \code{int}s, they find
the tightest at the powers of the tests \file{elementary} and
\file{ieee1788} there (odd and even exponents, negative bases, exponents of
64 and 128 bits), the tightest at 5 powers beyond $2^{53}$ where $n$ is a
double, and within three doubles at 4 where it is none, of intervals and of
expressions, against mpmath. They check that $x^2$ is \code{sqr(x)} over
400 intervals near $1$, $3$, $0.75$ and $2^{-500}$, and bit for bit, the
sign of a zero included, in each of its forms (\code{pow(x, 2)},
\code{2u}, \code{2.0}, \code{[2]}, \code{gaol\_ieee1788::pown(x, 2)} and
\code{pow(x, [2])}), over the special bounds and 90\,000 random intervals,
in each direction of the calling code and with each flush mode; that
$(-c)^n$ is $\pm c^n$, and that the bound from a corner is the same whether
the other corner is computed with it or alone, at 1\,088 points near $1$,
$1.5$ and $3$ and at 12\,736 more within 16 doubles of $1$. ITF1788, the
test suite of Chapter~\ref{chap:computing}, gets the tightest in its 163
cases of \ieee{pown}.

\paragraph{The real power.}
\code{pow(x, y)} and \code{pow(x, d)} for a non-integer $y$
(\code{gaol\_pow\_hybrid()}, \code{gaol\_pow\_real()}), and
\code{gaol\_ieee1788::pow(x, y)} for any $y$, compute \ieee{pow} on
$x\cap\itv{0}{+\infty}$: CORE-MATH's \code{pow} at the corners of the box
$x\times y$ where $x^y$ is the least and the greatest, which the places of
the bounds about $1$ and $0$ give, correctly rounded upward. The upper bound
is the value at the upper corner, and the lower bound is the double below
the value at the lower corner, unless $x^y$ is a double there. It is one
for $x=2^p$ and an integer $p\,y$ of $[-1074, 1023]$, and for
$x=c^{2^k}\cdot2^{f\cdot2^k}$ ($c\geq3$ odd) and $y=a/2^k$ ($a>0$, odd if
$k>0$) where $c^a<2^{53}$ and $c^a\cdot2^{f\cdot a}$ is within the doubles,
$y$ being then a multiple of $1/32$ in $(0, 33]$; the lower bound is then
$x^y$ itself, which a test on $x$ and $y$ proves with integers and exact
operations, leaving at once for nearly every $y$, where calling CORE-MATH
downward too would double the cost of the power.\newinvfive{} The result
is $1$ where a corner has the base $1$ or the exponent $0$. At a corner
where the base is $0$ or $+\infty$ or the exponent infinite, it is the
limit of $x^y$ there, $0$ or $+\infty$, which is the bound itself; GAOL~4
computed $\exp(y\log x)$ there, which gave bounds one to hundreds of
doubles beyond.\newinvfive{} $0^y$ is $0$ for $y>0$, and $\emptyset$ for
$y\leq0$. An integer $y$ beyond the \code{int}s, in
\code{gaol\_ieee1788::pow}, takes the corners at the bounds of $x$, $x^y$
being monotone in $x$, where the \code{pow} of \code{gaol} takes
\ieee{pown}; GAOL~4 gave $\itv{-\infty}{+\infty}$.\newinvfive

The result is tightest, at the corners and at their limits, those at the
bounds of $x$ for an integer $y$ beyond the \code{int}s included. For a
degenerate integer $y$ within the \code{int}s, it is the integer power
above, one double beyond the tightest where $x^y$ is within $2K$ of a
double.

\begin{Ex}
\code{pow([4], 0.5)} is $\itv{2}{2}$ and \code{pow([4, +oo], 0.5)} is
$\itv{2}{+\infty}$. \code{gaol\_ieee1788::pow([0x1.fffffffffffffp+1], [7])}
has its lower bound, \code{0x1.ffffffffffff8p+13}, one double below the
tightest, \code{0x1.ffffffffffff9p+13}: the integer power is within $2K$ of
a double there.
\end{Ex}

The tests find the bounds within 1 double at points and over 385 boxes; the
tightest, or one double beyond for one integer power, on the 117 boxes of
the test \file{ieee1788}, the limits at $0$ and $+\infty$ included, against
mpmath; the tightest on the 1\,347 cases of \ieee{pow} of ITF1788; the
tightest at 5 integer exponents beyond the \code{int}s; and the tightest,
both bounds, at the corners of 84\,070 boxes of a base and an exponent
(powers of two, perfect powers, midpoints, at the edges of the doubles too,
and their neighbours), against $x^y$ computed exactly at the 8\,981 where
it is rational, 3\,211 of them with a power that is a double, and compared
exactly at the 12\,540 where it is the $2^k$-th root of a rational,
$k\leq6$.

\paragraph{Exponentials and logarithms.}
\code{exp}, \code{exp2}, \code{exp10}, \code{log}, \code{log2} and
\code{log10} are CORE-MATH's functions at the bounds, correctly rounded
upward: the value at the right bound, and the double below the value at the
left bound, unless that value is a double. The logarithms work on
$x\cap\itv{0}{+\infty}$, give $\emptyset$ when $x$ holds no positive
number, and $\log(0)=-\infty$; the exponentials have a lower bound at least
$0$. The values that are doubles are given exactly: $\exp(0)=1$ and
$\log(1)=0$; for \code{exp2}, \code{exp10}, \code{log2} and
\code{log10},\newinvfive{} $2^x$ for an integer $x$ of $[-1074, 1023]$ and
$10^x$ for an integer $x$ of $[0, 22]$ ($10^{23}$ is no double, and no
negative power of ten is), $\log_2 x$ for a power of two and $\log_{10}x$
for a power of ten of $[10^0, 10^{22}]$. All six are tightest. The tests
find \code{exp} tightest at about 9\,000 doubles of every magnitude, \code{log} tightest at the 75
doubles of the tests and the million intervals of the comparison with the
other libraries (\pxref{sec:comparison}), and the other four tightest over
6\,483 doubles each, and at the values that are doubles.

\subsection{Trigonometric and hyperbolic functions}\label{sec:accuracy-trig}

These functions must be valid, and the standard recommends them to be
accurate. Every one of them is CORE-MATH's function at the bounds of the
pieces where it is monotonic, so each is tightest where the function is
evaluated at the bounds; the periodic functions ask first on which piece
the interval lies. Every elementary function of this section and of the
one before is that of CORE-MATH, whose sources are in \file{3rd/math-core}
(\pxref{sec:math-libraries}).

\paragraph{Cosine and sine.}
For \code{cos(x)}, the bounds divided by an enclosure of $\pi$, rounded
outward, tell whether $x$ lies on one piece where the cosine is monotonic;
rounded inward, they tell whether an extremum is within $x$ for sure. Where
they cannot tell, a bound being within about $|x|\cdot2^{-52}$ of a
multiple of $\pi$, or beyond $2^{52}$, the signs of CORE-MATH's \code{sin}
at the bounds do, exactly, $x$ being narrower than $2\pi$ (and the sign at
the middle of $x$ from $\pi$ on): signs that differ show one extremum, $-1$
or $1$ according to their order. Then CORE-MATH's \code{cos} is taken at
the bounds, correctly rounded upward (the double below the value at the
bound where the cosine is the least), within $\itv{-1}{1}$, and $\cos(0)=1$ exactly. \code{sin(x)}
proceeds as \code{cos}, the bounds divided by an enclosure of $\pi$ minus
$1/2$, the signs of CORE-MATH's \code{cos} where the quotients cannot tell,
and $\sin(0)=0$ exactly.

Both are tightest where the function is evaluated at the bounds, at every
magnitude, and give $-1$ and $1$ exactly at the extrema. They give
$\itv{-1}{1}$ from the width $2\underline{\pi}$ on, $\underline{\pi}$ being
the double below $\pi$, which is then the tightest result: an interval that wide misses an arc shorter than $2.5\cdot10^{-16}$, where the sine and the cosine stay within $10^{-31}$ of $\pm1$. That the test of a width below
$\underline{\pi}$ gives the tightest bounds is measured, not proved: the
tests find the tightest bounds of \code{sin} and \code{cos} over 29\,400
intervals of width about $\pi$ whose bounds are next to extrema of the
sine. They find them tightest for all doubles, and over intervals next to
the extrema and of width about $\pi$ and $2\pi$.

\paragraph{Tangent.}
\code{tan(x)} is $\itv{-\infty}{+\infty}$ when $x$ is wider than
$\underline{\pi}$ or holds a pole: $(x+[\pi/2])/[\pi]$ at the bounds,
rounded outward, tells that no pole is within $x$ when the integer parts
agree, and rounded inward that one is for sure; where they cannot tell, the
signs of CORE-MATH's \code{cos} at the bounds do, exactly, a pole being
within $x$ when they differ. Otherwise the result is CORE-MATH's \code{tan}
at the bounds, correctly rounded upward, and $\tan(0)=0$ exactly. It is
tightest between two poles, at every magnitude, and
$\itv{-\infty}{+\infty}$ exactly when a pole is within $x$, but for the
widths above $\underline{\pi}$ and below $\pi$, a theoretical case. The
tests find it tightest for all doubles, and over intervals next to the
poles, the interval from \code{-M\_PI\_2} to \code{M\_PI\_2} and the
intervals of width $\underline{\pi}$ or next to it between two poles
included.

\paragraph{Inverse trigonometric functions.}
\code{asin(x)} and \code{acos(x)} are CORE-MATH's functions at the bounds
of $x\cap\itv{-1}{1}$, correctly rounded upward; at $-1$, $0$ and $1$ they
give $0$ exactly or the tightest bounds of $\pm\pi/2$ and $\pi$.
\code{atan(x)} is CORE-MATH's \code{atan} at the bounds; at $0$, $\pm1$
and $\pm\infty$ it gives $0$ exactly or the tightest bounds of $\pm\pi/4$
and $\pm\pi/2$. \code{atan2(y, x)} is CORE-MATH's \code{atan2} at the
corners of the box $y\times x$ where the angle is the least and the
greatest, which the signs of the bounds give, correctly rounded upward; it
gives $0$, $\pm\pi/4$, $\pm\pi/2$ and $\pi$ where a corner is on an axis,
on a diagonal of the right half-plane or at an infinity, $\itv{-\pi}{\pi}$
when the box has points on the half-line $y=0$, $x<0$, where the angle is
$\pi$, and points below it, and $\emptyset$ for the box $\{(0, 0)\}$. The
four are tightest, and exactly $0$ or the tightest enclosures of the
multiples of $\pi/4$ named where a bound is one of them. The tests check
them within 1 double, and tightest at the special points.

\paragraph{Hyperbolic functions.}
\code{sinh(x)}, \code{tanh(x)} and \code{asinh(x)} are CORE-MATH's
functions at the bounds, correctly rounded upward, \code{tanh} within
$\itv{-1}{1}$; they give $0$ at $0$ exactly, $\sinh(x)$ beyond the largest
double for $|x|\geq711$, and $\tanh(x)$ beyond $\pm(1-2^{-53})$ for
$|x|\geq20$, where $1-|\tanh(x)|<2^{-54}$. \code{cosh(x)} takes CORE-MATH's \code{cosh} at the bound of largest magnitude for the upper bound, and at the bound of smallest magnitude for the lower one, which is $1$ when $x$ contains $0$, and a value beyond the largest double for
$|x|\geq711$. \code{acosh(x)} and \code{atanh(x)} proceed as \code{sinh},
on $x\cap\itv{1}{+\infty}$ and $x\cap\itv{-1}{1}$ respectively, with
$\mathrm{acosh}(1)=0$ and $\mathrm{atanh}(0)=0$ exactly. The hyperbolic arc
tangent is defined on $(-1, 1)$, $-\infty$ and $+\infty$ being its limits
at $-1$ and $1$, which are bounds but no values: the result is $\emptyset$
when $x$ holds no point of $(-1, 1)$, as for $\itv{1}{1}$, $\itv{1}{5}$,
$\itv{-1}{-1}$ and $\itv{-5}{-1}$.\newinvfive{} The six are tightest, and
exact at the points given. The tests check them within 1 double, tightest
at $0$, at $1$, and beyond $711$ and $20$, and $\emptyset$ for those four
intervals.

\subsection{Recommended functions}\label{sec:accuracy-recommended}

The recommended functions of Table~10.5 of the standard must be valid, and
the standard recommends them to be accurate. All of them are tightest but
the roots of an order other than 1, 2 and 3, which are within 2 doubles of
the tightest, and
\code{compoundm1}, which is tightest but in a case out of $10^{17}$ or so.

\paragraph{Roots.}\findex{nth\_root}%
\code{nth\_root(x, n)}, for $n>0$, computes \ieee{rootn}. For $n=1$ it is
$x$, and for $n=2$ it is \code{sqrt}. For $n=3$ it is CORE-MATH's
\code{cbrt} at the bounds, correctly rounded upward on their magnitudes,
the root of a negative number being the opposite of the root of its
magnitude: the tightest bounds, and exact where the cube root is a double,
which cubing the value tells; it takes 50~ns rather than 189 with the
search below, on an Intel i7-1185G7 with Clang 18.1.\newinvfive{} For the
other orders, on $x$ (odd $n$, the root of a negative number being again
the opposite of the root of its magnitude) or on $x\cap\itv{0}{+\infty}$
(even $n$), the lower bound is the largest double $l$ with $l^n$ rounded
upward at most the bound of $x$, and the upper bound the smallest double
$u$ with $u^n$ rounded downward at least it, which proves them. They are
looked for from CORE-MATH's \code{pow} with the exponent $1/n$ rounded,
after a step of Newton's method, by steps that double then by bisection,
two powers where the start is next to the root; the roots of $0$ and $1$
are $0$ and $1$. The orders 1, 2 and 3 are tightest, and so is any order
where the bound of $x$ is an $n$-th power of a double. Otherwise the bounds are the tightest or up to 2 doubles beyond, whatever the magnitude of $x$: the $n-1$ rounded products of $l^n$ move the root by less than $2^{-52}$ relatively. Two doubles beyond is not accurate where the order is large: $\code{nth\_root([0x1.16ef146ba87dap-473], 1000)}$ has its upper bound two doubles above the tightest one, one above the accurate one. An order beyond the \code{int}s and
the \code{unsigned int}s takes CORE-MATH's \code{pow} at the corners, with
the exponent $1/n$ as an interval, within 1 double of the
tightest.\newinvfive{} The tests find the roots within 2 doubles for all
doubles, 1 double being the largest distance found, among them the $x$
below $2^{-1020}$, subnormal ones included, with $n$ up to $2^{32}-1$,
against exact arithmetic and mpmath, and the tightest on the values of the
test \file{ieee1788} beyond the \code{unsigned int}s.

\code{nth\_root(x, q)}, for $q<0$, is $1/x^{1/|q|}$, the inverse of the
root above: the domain is then $\RSet\setminus\{0\}$ for an odd $q$ and
$(0, +\infty)$ for an even one, which taking the inverse gives, an interval
holding $0$ having for inverse the hull of the values away from it;
\code{nth\_root(x, 0)} is $\emptyset$.\newinvfive{} It is valid, and within 4 doubles of the tightest bounds: the root above, then the tightest division. It is not always accurate: \code{nth\_root([2.25], -127)} has its upper bound 4 doubles above the tightest one, 3 above the accurate one. The tests find that it encloses the root over 1.5 million
values against a reference in \code{long double}, and that it is the
inverse of the positive root over 8\,000 values.

\paragraph{Exponentials minus one and logarithms of one plus.}
\code{expm1(x)}, \code{exp2m1(x)} and \code{exp10m1(x)} are CORE-MATH's
functions at the bounds, correctly rounded upward, the double below the
value at the left bound unless that value is a double, which $b^x-1$ is at
$0$, and at the integers of $[-53, 53]$ for $b=2$ and of $[0, 15]$ for
$b=10$; the result is within $\itv{-1}{+\infty}$.\newinvfive{}
\code{log1p(x)}, \code{log2p1(x)} and \code{log10p1(x)} are CORE-MATH's
functions at the bounds of $x\cap\itv{-1}{+\infty}$, the value at $-1$
being the limit $-\infty$, and $x$ holding no number above $-1$ giving
$\emptyset$; the double below the value at the left bound unless it is a
double, which $\log_b(1+x)$ is at $0$, at $2^k-1$ for $b=2$ ($k$ in
$[-53, 53]$) and at $10^k-1$ for $b=10$ ($k$ in $[0, 15]$)
only.\newinvfive{} The six are tightest. The tests find the three
exponentials tightest over 400\,000 points and intervals each, and the
three logarithms tightest over 40\,000 intervals each, their bounds drawn
among those points and their neighbours.

\paragraph{The functions of $\pi x$.}
For \code{sinpi(x)} and \code{cospi(x)}, the extrema are at the multiples
of $1/2$, which are doubles: doubling the bounds, which is exact, and
taking their ceiling and floor give exactly the integers $t=2x$ within $x$,
and $\sin(\pi x)$ has its maximum $1$ at the $t$ congruent to $1$ modulo
$4$ and its minimum $-1$ at those congruent to $3$ (the cosine at $0$ and
$2$); otherwise the result is CORE-MATH's function at the bounds, exact at
the multiples of $1/2$. Beyond $2^{61}$ two distinct doubles are whole
periods apart, and the result is $\itv{-1}{1}$.\newinvfive{} Both are
tightest, at every magnitude: $\mathrm{sinpi}(10^{17})$ is $0$, where
\code{sin} of $\pi\cdot x$ gave $\itv{-1}{1}$. \code{tanpi(x)} proceeds as
\code{sinpi}, its poles being at the odd $t$: $\itv{-\infty}{+\infty}$
with a pole within $x$, $-\infty$ or $+\infty$ for a pole at a bound, and
$\emptyset$ for a pole alone, where the function has no
value.\newinvfive{} It is tightest. The tests find the three tightest over
400\,000 points and intervals each.

\code{atanpi(x)}, \code{asinpi(x)} and \code{acospi(x)} are CORE-MATH's
functions at the bounds, \code{asinpi} and \code{acospi} on
$x\cap\itv{-1}{1}$, \code{acospi} decreasing; they are exact at $0$ and
$\pm1$ (and at $\pm\infty$ for \code{atanpi}), the only doubles where the
values are doubles (\code{asinpi}$(\pm1/2)=\pm1/6$, \code{acospi}$(1/2)=1/3$ and \code{acospi}$(-1/2)=2/3$ are rational too, but no doubles).\newinvfive{} \code{atan2pi(y, x)} is the analysis of
the box of \code{atan2}, the angles divided by $\pi$: $\itv{-1}{1}$ for a
box crossing the half-line $y=0$, $x<0$, and the corners elsewhere. The
values at the multiples of $1/4$ ($0$, $\pm1/4$, $\pm1/2$, $\pm3/4$, $1$,
and $-1$ as a limit) are exact, and CORE-MATH's \code{atan2pi} gives the
others, which are irrational.\newinvfive{} The four are tightest. The
tests find them so over 400\,000 points and intervals each (40\,000
intervals for \code{asinpi}, and the arguments next to $\pm1$ where the accurate phase of
CORE-MATH's \code{asinpi} shifted a 64-bit integer by 65 to 69 bits, which \file{3rd/README.md} describes),
and \code{atan2pi} over 200\,000 boxes.

\paragraph{Hypotenuse and inverse square root.}
\code{hypot(x, y)} is CORE-MATH's \code{hypot} at the point of the box
nearest to the origin, $(\mathrm{mig}\,x, \mathrm{mig}\,y)$, and at the
farthest, $(\mathrm{mag}\,x, \mathrm{mag}\,y)$. The lower bound is the
value itself where it is a double: writing $a$, $b$ and $h=\sqrt{a^2+b^2}$
as odd integers times powers of two, $a^2+b^2=h^2$ can only hold where $h$
has the power of two of the one of $a$, $b$ whose power is the smaller, and $h^2-b^2=a^2\cdot4^d$
is then compared exactly on 107 bits.\newinvfive{} \code{rsqrt(x)} is
CORE-MATH's \code{rsqrt} at the bounds of $x\cap\itv{0}{+\infty}$,
decreasing, the value at $0$ being the limit $+\infty$, and $x$ holding no
positive number giving $\emptyset$; it is exact at the powers of
$4$.\newinvfive{} Both are tightest. The tests find \code{hypot} tightest
over 200\,000 boxes, among them Pythagorean triples scaled by powers of two
and the doubles next to them, and \code{rsqrt} over 40\,000 intervals.

\paragraph{Compound minus one.}\findex{compoundm1}%
\code{compoundm1(x, y)} computes $(1+x)^y-1$ with the analysis of the box
of \code{pow} for the base $1+x$: the value at the corner each bound is
at, since $(1+x)^y-1$ increases with $y$ for $x>0$ and decreases for
$x<0$, increases with $x$ for $y>0$ and decreases for $y<0$, and is $0$ on
$x=0$ and on $y=0$. The limits $-1$ and $+\infty$ are the bounds at
$x=-1$, at $x=+\infty$ and for an infinite $y$, and the empty set is the
result for $x<-1$ and for $x=-1$ with $y\leq0$, where the function has no
value. CORE-MATH has no \code{compoundm1}, and
$\exp(y\,\mathrm{log1p}(x))-1$ computed with doubles multiplies the
relative error of $\mathrm{log1p}$ by $|y\,\mathrm{log1p}(x)|$: up to
1\,100 doubles from the value near the overflow, over 300 draws with the
\code{expm1} and \code{log1p} of Python. The value at a corner is
therefore computed in \file{gaol/math/gaol\_compoundm1.h} with integers
only. A fast phase in 64-bit fixed point comes first, whose error bound is
proved (27 to 200 units of $2^{-63}$ relatively, 4.6 to 10 times the worst
error found); it settles 95\,\% of the pairs for $|y\,\mathrm{log1p}(x)|$
below $0.35$, 92\,\% from $0.35$ to $1$, 90\,\% from $1$ to $2$, 83\,\%
from $2$ to $4$ and 71\,\% from $4$ to $8$, and none past, nor for
$|\mathrm{log1p}(x)|$ above $8$, $x<-0.99966$ or $x>2\,980$. The other
pairs go on to a 128-bit significand in balls, which add the error of each
operation to a bound of the error. Both phases start from the
$\mathrm{log1p}(x)$ of CORE-MATH, corrected by one step of Newton's method
($d=(1+x)\exp(-y_0)-1$, then $y_0+d-d^2/2+d^3/3$), multiply it by $y$,
then take $\exp-1$ by a series of 27 terms after the reduction by a
multiple of $\log2$. The ball is no wider than $2^{-111}$ relatively for
$|y\,\mathrm{log1p}(x)|$ below $746$ ($2^{-118}$ below $10$, $2^{-120}$
below $1$), and the bounds are the doubles around it. A corner takes
$0.1\,\mu$s where the fast phase settles it and $0.6$ to $0.9\,\mu$s in
the balls, on an Intel i7-1185G7 with GCC 15: $0.12\,\mu$s for a point and
$0.23\,\mu$s for a box of $x$ between $-0.3$ and $0.5$ and $y$ between
$0.1$ and $3$, instead of $2\,\mu$s with the balls alone.

The cases where the value is a double are told apart. For $y=-1$ the value
is the quotient $-x/(1+x)$ of two integers, which is divided exactly, and
the value is the double where the division is exact. The value is the
double $d$ where $(1+x)^y=1+d$, which integers prove for $y=P/2^k$:
$s=1+x$ a power of two with $y\log_2 s$ an integer, or $s$ the $2^k$-th
power of an integer $R$ times a power of two with $R^P=1+d$
($2.375=1.5^3-1$, $1.25$, $0.5=\sqrt{2.25}-1$, $2^m-1$ for $|m|\leq53$);
it is $y\,x$ on the side $y(y-1)$ tells where $y\,x$ is a double and
$x(|y|+2)\leq1/2$, the ball not telling the side where $x\,|y-1|$ is below
$2^{-119}$; and $2^m-1$ is between the double below $2^m$ and $2^m$ for
$m$ past 53, which the ball cannot tell from $m=118$ on.\newinvfive

The function is tightest, but where the value is within $2^{-111}$
relatively of a double that is not the value, none of the cases above
tells it and the fast phase, which tells apart the doubles that are
$2^{-55}$ to $2^{-58}$ apart relatively or more, has not settled the pair:
about one pair $(x, y)$ in $10^{17}$ drawn at random for the balls, 5\,\%
of the pairs going to them, and none among the 5\,854 pairs of the tests;
the result is then one double wider on that side. The contrived values
$y=-1$ and $x=\pm(2^e+2^{2e-1})$ for $e$ from $-51$ to $-40$, where
$1/(1+x)=1-x+x^2-x^3\ldots$ has its cubic term cancelled and is within
$2^{-155}$ of a double, are the tightest ones now, with the division. The
tests find the function tightest over 5\,854 pairs $(x, y)$ of every
magnitude, among them the values that are doubles, the tiny $x$ down to
the subnormals, $y=-1$ (719 values of $x$, the contrived ones among them,
633 of them divided), the edges of the fast phase
($|\mathrm{log1p}(x)|=1/2$ and $8$, $y\,\mathrm{log1p}(x)=8$ and
$\log(2)/2$), the thresholds of the reduction and of the cuts, and 850
boxes (the definition read from the hull of the corners), against mpmath at
1\,500 to 7\,000 bits and fractions; the balls of $\mathrm{log1p}$ and
$\exp-1$ are checked against mpmath to 256 bits.

\subsection{The error functions}\label{sec:accuracy-erf}

IEEE 1788-2015 does not name \code{erf} and \code{erfc}. GAOL computes them
with CORE-MATH's \code{erf} and \code{erfc} at the bounds, correctly
rounded upward, \code{erf} being increasing and \code{erfc} decreasing: the
value at one bound, and the double below the value at the other one unless
that value is a double, which it is at $0$ ($\mathrm{erf}(0)=0$,
$\mathrm{erfc}(0)=1$) and at $\pm\infty$ (the limits $\pm1$ of
$\mathrm{erf}$, $2$ and $0$ of $\mathrm{erfc}$); the results are within
$\itv{-1}{1}$ and $\itv{0}{2}$. $\mathrm{erf}(x)$ is within $2^{-53}$ of
$\pm1$ from $|x|=5.87$ on, where its bounds are $1-2^{-53}$ and $1$, or
$-1$ and $-(1-2^{-53})$; $\mathrm{erfc}(x)$ is subnormal from $x=26.55$
on, and below $2^{-1074}$ from $x=27.22$ on, where its lower bound is $0$.
Neither function is known to take a value that is a double at any other
double, and the sources of CORE-MATH treat no such case.\newinvfive{} Both
are tightest. The tests find them so at the 409 doubles of
\file{elementary\_values.h} (tiny arguments, changes of binade, the
arguments where the values round to $\pm1$, $2$ or $0$ or become subnormal,
and the hard cases of CORE-MATH's \file{erfc.c}) and over the intervals
between them, against mpmath, and over 9\,225 doubles and 80\,000 intervals
against CORE-MATH called in both directed roundings.

\subsection{Reverse functions}\label{sec:accuracy-reverse}
\cindex{reverse functions!accuracy}%

The relational functions of GAOL compute \code{f\_rel(J, I)}, the hull of
the $x$ of $I$ whose image is in $J$: they are the reverse functions of
Table~10.1 of the standard with the arguments in another order
(\pxref{sec:relational-arithmetic}). Each computes the preimage of $J$
with the functions above, rounded outward, and intersects it with $I$: it
is valid, which the standard requires. The test \file{tests/reverse.cpp}
checks, on the cases of the minimal tests of libieeep1788 and on random
ones, that they are also accurate, which the standard recommends for
inf-sup types, the hulls being computed from the definition of each
function, apart from GAOL (\pxref{sec:test-programs}); the test
\file{tests/other\_functions.cpp} checks that they keep the values they
are given, within the numbers of doubles below.

\begin{itemize}
\item[$\bullet$] \textbf{Square root.} \code{sqrt\_rel(c, x)}
  (\ieee{sqrRev}) is accurate; the tests find it accurate, and within 1
  double of the value it keeps.
\item[$\bullet$] \textbf{Absolute value.} \code{invabs\_rel(c, x)}
  (\ieee{absRev}) is tightest: its bounds are those of $J$, of $-J$ and of
  $I$.
\item[$\bullet$] \textbf{Integer power, $p>0$.}
  \code{nth\_root\_rel(c, p, x)} (\ieee{pownRev}) is accurate, its bounds
  being the roots of \code{nth\_root()}; beyond the \code{unsigned int}s it
  is the hull of the parts of $x$ in those roots and, for an even $p$, in
  their opposites.\newinvfive{} The tests find it accurate, within 1 double
  of the tightest, and within 2 doubles of the value it keeps.
\item[$\bullet$] \textbf{Integer power, $p\leq0$.}\newinvfive{} For $p=0$,
  \code{nth\_root\_rel(c, 0, x)} is $x$ where $1$ is in $c$, and
  $\emptyset$ elsewhere: exact. For $p<0$, $c$ is cut at $0$, and the part
  at least $0$ and, for an odd $p$, the part at most $0$ each give the
  roots \code{nth\_root(part, p)}, the root then the tightest division, up
  to 4 doubles beyond the tightest bounds; the result is the hull of their
  parts in $x$, $0$ left out. The hull of the roots of $c$, which joins its
  parts, would keep the part of $x$ between them, and the roots reach $0$
  where $c$ is unbounded, $0$ having no power of negative exponent. The
  result is valid, and within 2 doubles of the accurate bounds of 12.10.1, one
  more than the trigonometric functions: the tests find it $\emptyset$ where
  the hull is $\emptyset$ in the cases written out, and of the 18\,052 random
  cases with $p<0$ of
  \file{tests/reverse\_values.py}, 12\,495 have a nonempty result with
  nonempty accurate bounds, 521 of them one double beyond, 25 two, none
  more. It is at most 2 doubles from the tightest bounds on the table of
  \file{tests/reverse.cpp}, 3 on those random cases (60 of them), and 4 at
  most, found for $c=\itv{2.25}{2.25}$ and $p=-127$. Where $p=-1$ and
  $1/c$ reaches beyond the largest double, $c$ being next to $2^{-1074}$, a
  bound saturates at $\pm$\code{DBL\_MAX}, as \code{div\_rel} does: valid,
  but beyond 4 doubles of the tightest, and not empty where the hull is:
  the result for $c=\itv{2^{-1074}}{2^{-1074}}$, $p=-1$ and
  $x=\itv{0}{\code{DBL\_MAX}}$ is $\itv{\code{DBL\_MAX}}{\code{DBL\_MAX}}$.
\item[$\bullet$] \textbf{Trigonometric functions.} \code{asin\_rel},
  \code{acos\_rel} and \code{atan\_rel} (\ieee{sinRev}, \ieee{cosRev},
  \ieee{tanRev}) are accurate, or one double beyond: the pieces of the
  preimage are $k\pi$, enclosed within about one double from $\pi$ in
  double-double, plus or minus the inverse function of $J$, which is added
  before the sum is rounded; beyond $2^{52}$ a bound of $x$ is kept, and an
  $x$ of a single double is decided by the image of the function. The tests
  find them accurate but for one double, with the bounds of $x$ beyond
  $2^{52}$, and within 6 doubles (4, 5 and 3 found) of the value they keep
  from $1$ to $2^{50}$, away from the points where the inverse function
  magnifies the width of the image.
\item[$\bullet$] \textbf{Hyperbolic functions.} \code{acosh\_rel(c, x)}
  (\ieee{coshRev}) is accurate; the tests find it within 16 doubles of the
  value it keeps. \code{asinh\_rel} and \code{atanh\_rel}, named after
  \ieee{coshRev} (\ieee{sinhRev} and \ieee{tanhRev} are not in Table~10.1),
  are accurate; \code{atanh\_rel(J, x)} is $\emptyset$ where $J$ holds no
  point of $(-1, 1)$, the values of the hyperbolic tangent.\newinvfive{}
  The tests find them within 10 and 26 doubles of the value they keep, and
  $\emptyset$ for $J=\itv{1}{1}$ and $\itv{-1}{-1}$.
\item[$\bullet$] \textbf{Division.} \code{div\_rel(c, b, x)}, and
  \code{c \% b} for $x=\itv{-\infty}{+\infty}$ (\ieee{mulRev}), are
  accurate, and tightest where the bounds of the quotients are doubles. The
  tests find them accurate, the tightest over the cases of \code{\%}, and
  \code{div\_rel} within 2 doubles of the value it keeps.
\item[$\bullet$] \textbf{Powers.}\newinvfive{} \code{pow\_rel(c, b, x)} and
  \code{pow\_exponent\_rel(c, a, x)} (\ieee{powRev1}, \ieee{powRev2}) are
  tightest. The pieces of the solutions over the signs of the other
  argument ($y>0$, $y<0$; $x>1$, $x<1$), with $0^y=0$ for $y>0$, $x^0=1$
  and $1^y=1$, end at $c^{1/y}$ or $\log(c)/\log(x)$, $c$, $x$ and $y$
  being bounds of the arguments; their bounds are the largest double $d$
  with $d^y$ (or $x^d$) proved by CORE-MATH's \code{pow}, correctly
  rounded, to be on one side of $c$, and the smallest one on the other
  side. A bound of $x$ where the solutions start between it and the next
  double, and whose power is within one double of $c$, is kept, which only
  knowing whether that power is a double would settle: the result is then
  the tightest one for $x$ one double wider. The tests find them tightest
  on the 804 cases of ITF1788 and on 1\,244 others, random or whose
  solutions are doubles, within the tightest result for $x$ one double
  wider on 966 cases where a bound of $x$ is next to an end of the
  solutions, and keeping the 54\,305 points of 6\,000 random boxes whose
  power is within $c$.
\item[$\bullet$] \textbf{Maximum, minimum, sign and floor.}\newinvfive{}
  \code{max\_rel(z, y, x)}, \code{min\_rel(z, y, x)},
  \code{sign\_rel(z, x)} and \code{floor\_rel(z, x)}, whose reverse
  functions Table~10.1 does not name, are tightest: their bounds are those
  of $x$ and $z$, $0$, $\lceil\leftBound{z}\rceil$ and
  $\lfloor\rightBound{z}\rfloor+1$, the open ends of the sets of
  \code{sign\_rel} ($0$) and \code{floor\_rel}
  ($\lfloor\rightBound{z}\rfloor+1$) included, the last rounded upward
  beyond $2^{53}$. The tests check them against the definitions on every
  interval of half-integer bounds, on cases written out and by sampling,
  and \code{floor\_rel} at every magnitude, against
  $\lfloor\rightBound{z}\rfloor+1$ computed exactly.
\item[$\bullet$] \textbf{Two-argument arc tangent.}\newinvfive{}
  \code{atan2\_rel\_y(c, b, x)} and \code{atan2\_rel\_x(c, a, x)}
  (\ieee{atan2Rev1}, \ieee{atan2Rev2}) are accurate. The angles of $c$ in
  $(-\pi, \pi]$ are cut at the axes into four cones, whose sides have the
  slopes $\tan\theta$ at the bounds $\theta$ of $c$ (CORE-MATH, correctly
  rounded, the slope of a subnormal $\theta$ being $\theta$ itself), and
  $0$ or $+\infty$ at the axes; the bounds of the box, and their products
  and quotients with the slopes, are rounded outward. The result is within
  2 doubles of the tightest where both are nonempty; a box within one
  rounding of the cone may give a point that the tightest enclosure, the
  empty set, does not hold. The tests find them within 2 doubles of the
  tightest over 1\,969 cases (1\,263 nonempty) against mpmath at 3\,000
  bits, around the multiples of $\pi/4$, the axes, the half-line $y=0$,
  $x<0$, and with zero, subnormal and infinite bounds, and keeping every
  point whose angle is in $c$, over about 600\,000 such points of the
  boxes, among them points of their sides next to the sides of the cones.
\end{itemize}

\begin{Rem}\label{rem:accuracy-reverse-point}
Where a bound of \var{J} is one double from the image of a bound of the
preimage, the reverse functions keep a point of \var{I} whose image is
just outside \var{J}:
\code{sqrt\_rel(\itv{1+2^{-52}}{+\infty}, \itv{-1}{1+2^{-52}})} is
\itv{-1}{1+2^{-52}}, whose tightest enclosure is \itv{1}{1+2^{-52}}. They
are accurate all the same, the widened \var{J} of nextOut holding that
image.
\end{Rem}

\medskip
\ieee{mulRevToPair(b, c)} (10.5.5) is the pair of the two parts, each a
\code{div\_rel} of $c$ by a part of $b$ over $\itv{-\infty}{+\infty}$, the
part of $b$ at most $0$ and the part at least $0$:\newinvfive{} as tight as
\ieee{mulRev}, so the tightest where the bounds of the quotients are
doubles. The test \file{tests/ieee1788.cpp} checks it against exact
quotients, and the 172 cases of ITF1788 pass.

\ieee{cancelMinus(x, y)} and \ieee{cancelPlus(x, y)} (10.5.6) are
\code{cancel\_minus(x, y)} and \code{cancel\_plus(x, y)}, the tightest,
which 12.12.5 requires. Whether the result exists depends on $x$ being at
least as wide as $y$, which rounding the differences of the bounds cannot
decide when they are within a double of each other; each difference is
computed exactly, as the sum of its value rounded to nearest and of its
error, by the TwoSum of Knuth. The result is $\itv{-\infty}{+\infty}$
where there is no value, and $\emptyset$ for an empty $x$ and a bounded
$y$. They are checked over 200\,000 pairs against exact rational
arithmetic, 11\,524 of them of exactly equal widths.\newinvfive

\subsection{Set operations and numeric functions}\label{sec:accuracy-set}

The set operations, the constructors, the numeric functions and the
boolean functions are computed from the bounds, and most of them are
exact. The relations are described in Chapter~\ref{chap:relations}, the
reading and the writing of intervals in Chapter~\ref{chap:io}; here are
their accuracies.

\begin{itemize}
\item[$\bullet$] \textbf{Intersection and hull.} \code{x \& y} and
  \code{x | y} (\ieee{intersection}, \ieee{convexHull}) are taken on the
  bounds: exact.
\item[$\bullet$] \textbf{Constructors.} \code{interval(l, u)} and
  \code{interval(d)} (\ieee{numsToInterval}) store the bounds, and the
  empty set where $l>u$, $l=+\infty$, $u=-\infty$ or a bound is NaN:
  exact.
\item[$\bullet$] \textbf{Reading.} \code{gaol::textToInterval(s)} and
  \code{operator>{}>} (\ieee{textToInterval}) read the numbers, decimal,
  hexadecimal, octal or binary, compared exactly with the doubles around
  them; the rational numbers $p/q$ are divided as intervals; the uncertain
  form is read as its exact decimal bounds; the literals \code{[ ]},
  \code{[empty]}, \code{[entire]}, the infinite and the missing bounds, in
  any case of letters, are read (\pxref{sec:reading-intervals}). GAOL also
  reads expressions, such as \code{"sin(1)+0.1"}, computed with the
  operations above. The reading is tightest for the literals of IEEE
  1788-2015, $p$ and $q$ being doubles. The literal \code{[l, u]} is empty
  where $l$ is above $u$, numbers that are not doubles being compared
  exactly, but for \code{pi} and another number between the same two
  doubles, a number far beyond the doubles, or bounds that are computed,
  which give the interval from the lower bound of $l$ to the upper bound of
  $u$.\newinvfive{} It is valid for expressions. The tests find it
  tightest.
\item[$\bullet$] \textbf{Bounds.} \code{left()} and \code{right()}
  (\ieee{inf}, \ieee{sup}) return the stored bounds, and NaN for the empty
  set, where IEEE 1788-2015 has $+\infty$ and $-\infty$, and a zero lower bound of either sign, where \ieee{inf} gives $-0$ (12.12.8), as \code{gaol\_ieee1788::inf()} does: exact.
\item[$\bullet$] \textbf{Midpoint.} \code{midpoint()} (\ieee{mid}) is
  $(l+u)/2$ rounded to nearest, ties to even, or $l/2+u/2$ where $l+u$ may
  overflow, a bound being $2^{1023}$ or more in magnitude, which is tested
  first, so that $\itv{\code{DBL\_MAX}}{\code{DBL\_MAX}}$ raises no
  overflow exception;\newinvfive{} it is $0$, $-$\code{DBL\_MAX} and
  \code{DBL\_MAX} for the unbounded intervals, and NaN for the empty set:
  as 12.12.8 requires. \code{mid()} is the same midpoint rounded outward:
  tightest.
\item[$\bullet$] \textbf{Width and radius.} \code{width()} (\ieee{wid}) is
  $u-l$ rounded upward, NaN for the empty set; \code{rad()} and
  \code{mid\_rad(m, r)} (\ieee{rad}, \ieee{midRad}) give the greater of
  $m-l$ and $u-m$ rounded upward, $m$ being \code{midpoint()}, $+\infty$
  for the unbounded intervals and NaN for the empty set: as 12.12.8
  requires. \code{mag()}, \code{mig()} and \code{smig()} are taken on the
  bounds: exact. \code{hausdorff(x, y)} is the distance rounded upward:
  tightest.
\item[$\bullet$] \textbf{Relations.} \code{is\_empty()},
  \code{is\_a\_double()} and \code{set\_contains(d)} (\ieee{isEmpty},
  \ieee{isSingleton}, \ieee{isMember}), the set relations \code{set\_eq},
  \code{set\_contains} and \code{set\_leq}, \code{set\_strictly\_contains}
  and \code{set\_le}, \code{set\_disjoint} (\ieee{equal}, \ieee{subset},
  \ieee{interior}, \ieee{disjoint}, as Tables~10.3 and 10.4), the certainly
  relations \code{certainly\_leq}, \code{certainly\_le},
  \code{certainly\_geq} and \code{certainly\_ge} (\ieee{precedes},
  \ieee{strictPrecedes}, true when an interval is empty),
  \code{is\_entire()} and \code{is\_common\_interval()} (\ieee{isEntire},
  \ieee{isCommonInterval}, a common interval being nonempty and bounded,
  the NaN bounds of the empty set being no infinities\newinvfive) and
  \code{less(y)} and \code{strictly\_less(y)} (\ieee{less},
  \ieee{strictLess}, as Table~10.3, \ieee{strictLess} counting an infinite
  bound as below itself, with the empty cases of Table~10.4\newinvfive) are
  decided on the bounds: exact (Chapter~\ref{chap:relations}). The tests
  find them exact, and \code{less} and \code{strictly\_less} as Tables~10.3
  and 10.4 over 20\,000 pairs.
\item[$\bullet$] \textbf{Overlap.} \code{x.overlap(y)} and
  \code{gaol\_ieee1788::overlap(x, y)} (\ieee{overlap}) give the state of
  Table~10.7 among the 16 of \code{overlapping\_state}:
  \element{bothEmpty}, \element{firstEmpty} or \element{secondEmpty} for an
  empty operand, else the order of the four bounds, told as the integers of
  their bits (\code{gaol\_detail::bound\_key()}: $-0$ and $+0$ are the same
  bound, and an infinite bound is the same as itself); a point interval on
  a bound of the other one is \element{starts}, \element{finishes},
  \element{finishedBy}, \element{startedBy} or \element{equals}, never
  \element{meets} nor \element{metBy}.\newinvfive{} Exact. The tests read
  the set specification of Table~10.7 with quantifiers over a finite model
  of the real line, on 31\,449 pairs (the empty set, the 104 intervals of
  a grid of bounds, infinities, zeros of both signs, subnormals and
  neighbours of a double among them, and 20\,000 random pairs whose bounds
  meet); they check \ieee{disjoint}, \ieee{equal} and \ieee{subset} as
  masks of the states, the converse state when the operands are exchanged,
  the 48 cases of ITF1788 without decorations, and the same state with a
  mode that flushes the subnormals to zero.
\item[$\bullet$] \textbf{Writing.} \code{gaol\_ieee1788::intervalToText(x)}
  writes the bounds \code{[l, r]}, \code{[a]} for a point interval whose
  digits are exact (\code{[4]}, and \code{[0]} for a zero of either sign),
  or \code{[empty]}, in decimal with the digits of
  \code{interval::precision()} and a decimal point, whatever the global
  format and the locale.\newinvfive{} \code{operator<{}<} in the format
  \element{bounds}, and in \element{agreeing} for the intervals it writes
  with their bounds (those whose bounds share no nonzero leading digit, and those whose bounds are not both finite, nonzero and of the same sign), writes the same, with the
  digits and the flags of the stream, and a decimal point whatever its
  locale, so that the reader reads the text back; \element{agreeing}
  writes the other intervals with the digits their bounds share
  (\pxref{sec:output-format}). Each bound is written to nearest by the C
  library, compared exactly with the bound, and its last digit moved
  outward when it is on the wrong side. The format \element{hexa} writes
  the bounds in the hexadecimal-significand form of 13.4.1, \code{[a]} for
  a point interval and \code{[0x0p+0]} for a zero. The decimal formats are
  valid, and the tightest with the digits asked for where the C library
  rounds to nearest as it should; the hexadecimal format is exact. The
  tests find less than one unit of the last digit in decimal and the exact
  bits in hexadecimal, read the text of \ieee{intervalToText} back as a
  literal, over every format, precision and a comma locale, and read the
  text that \code{operator<{}<} writes for a point under that locale back
  as the point, or as an interval containing it.
\item[$\bullet$] \textbf{The formats for the eye.} \code{operator<{}<} in
  the formats \element{width} and \element{center} writes the
  \code{midpoint()} to nearest and the \code{rad()} rounded
  upward,\newinvfive{} as \code{c (+/- w)}, or \code{c} for a point
  interval, and \code{[empty]} for the empty set; the digits of $c$ are
  those of the precision, so that the text need not contain the interval.
  As doubles, $\itv{c-w}{c+w}$ contains $x$, $w$ being the smallest radius
  that does so, and a nonzero radius is never written $0$. The tests find
  the radius written at least the exact half-width and \code{rad()} rounded
  upward with less than one unit of its last digit, $\itv{c-w}{c+w}$
  containing $x$, and the midpoint written to nearest with at most half a
  unit of its last digit, over the special values, the subnormals, the
  largest doubles and random intervals, every precision and flag.
\item[$\bullet$] \textbf{Recovery.} \code{exact\_string(x)}, which
  \code{operator<{}<} writes with \element{hexa}, and
  \code{gaol::textToInterval(s)} (\ieee{intervalToExact},
  \ieee{exactToInterval}) meet the recovery requirement of 13.4: an
  interval written in hexadecimal and read again gives the same bounds, bit
  for bit, the empty set, the infinite bounds, a bound $-0$ and the
  subnormals included, but for a point interval of zero, written
  \code{[0x0p+0]} whatever the signs of its bounds and read back as
  $\itv{+0}{+0}$, the same set.\newinvfive{} Exact, over 2\,000 random
  intervals and the values written apart.
\end{itemize}

\section{The same bounds on every machine}\label{sec:accuracy-same-bounds}
%%------------------------------------------------------------------------
\cindex{accuracy!the same bounds on every machine}%
\cindex{zero!sign of a zero bound}%

GAOL is built with CORE-MATH only, whose sources are compiled into the
library (\pxref{sec:math-libraries}): the bounds of this chapter are the
bounds on every machine. They are therefore those of every build of GAOL,
and the test \file{tests/core\_math.cpp} checks them against CORE-MATH
called in the downward and the upward rounding, which are the tightest
bounds there are.

\medskip
The SSE2, the NEON and the FPU intervals (\pxref{sec:kernels}) compute the
integer powers with the same code, from exact products
(\code{pown\_products()} of \file{gaol/core/gaol\_interval\_core.cpp}), and
the square from one product per bound, and give the same bounds. The SSE2
intervals of GAOL~4 multiplied rounded products from the highest bit of
$n$, and their bounds and those of the FPU intervals differed for about half
of the intervals with a bound $0$ or infinite or whose power leaves the range
of the exact products.\newinvfive{} \code{pow(x, 2)} is \code{sqr(x)} in every form of
the integer power with the exponent 2 (\code{pow(x, 2)}, \code{2u},
\code{2.0} and \code{[2]}, \code{gaol\_ieee1788::pown(x, 2)}, and
\code{gaol\_ieee1788::pow(x, [2])} on the part of $x$ in
$\itv{0}{+\infty}$, which is $\itv{0}{0}$ for the part
$\itv{0}{0}$):\newinvfive{} the same function, so the same bits, the sign
of a zero bound being that of \code{sqr} in each build. In the earlier versions of GAOL v5, the integer power
made its square from exact products, which has the same bounds and gave a
zero bound of the other sign for part of the intervals holding $0$:
\code{pow(interval(0, 2), 2)} was $\itv{-0}{4}$ with the NEON intervals and
\code{pow(interval(-1, 2), 2)} was $\itv{-0}{4}$ with the FPU ones, where
\code{sqr} gives $\itv{0}{4}$; with the NEON intervals
\code{pow(interval::universe(), 2)} was $\itv{0}{+\infty}$ and is
$\itv{-0}{+\infty}$, as \code{sqr} has it. The exception flags are those
of \code{sqr} as well (\pxref{sec:fp-exceptions}): the corner products
that are not a bound of the result raise the inexact and the underflow
flags too (\code{pow([-1e-310, 1e154], 2)} raises underflow for the square
of its lower bound, which its lower bound $0$ leaves out), and the
overflow flag is raised when the upper bound is $+\infty$ only, the
smaller of two squares overflowing only with the larger: none of 14
million cases of each build had a finite upper bound with it.

\medskip
The bounds are the same, but the sign of a zero bound may differ between
the kernels, $+0$ and $-0$ being the same number.

\begin{Ex}
\code{sqr(interval(-2.0, 3.0))} is $\itv{-0}{9}$ with the SSE2 and the
NEON intervals, and $\itv{0}{9}$ with the FPU intervals.
\end{Ex}

So may the sign of a zero bound of a product or of a quotient. The SSE2
intervals take the larger of four products, and the quotient of the
corners their signs choose, with no branch, where GAOL~4 chose one product
by the signs of the bounds: \code{[-1, 0] * [-1, 0]} has the lower bound
$+0$ where it had $-0$.\newinvfive{} The FPU intervals take the larger of
four products too (for each bound, the lower one made of products with a
factor negated),\newinvfive{} but with a comparison and a selection that
the compiler compiles as it chooses, not an instruction: where two
products are zeros of opposite signs, the sign of the zero bound of a
product is the compiler's and the target's, $+0$ or $-0$. The lower bound
of \code{[-1, 0] * [-1, 0]} is $+0$ with GCC for 32-bit ARM and for arm64,
with Clang and with Visual C++ for x86, and $-0$ with GCC 13 for POWER9,
which compiles the selection into an \code{xsmaxcdp} that returns its
first operand for a $+0$ and a $-0$. The FPU intervals give the bits of
the SSE2 intervals, checked against a model of their kernels and not by
running them, for the products of finite bounds but for that sign, and for
the quotients by a divisor without zero: \code{[0, 1] / [-2, -1]} has the
upper bound $-0$, where the FPU intervals gave $+0$ before, the zero bound
of the dividend being divided. \code{div\_rel(K, J, I)} of the SSE2
intervals takes the code of the FPU and the NEON intervals where a bound of
$K$ is $2^{500}$ or more in magnitude, or a bound of $J$ that is not $0$
is below $2^{-500}$, the quotients of the vectors being liable to
overflow,\newinvfive{} so that a zero lower bound may have the sign of
that code: \code{div\_rel([0, 0x1.e2ab041732196p-5],
[2\textasciicircum-1074, 0x1.182p+9], [-DBL\_MAX, 0x1.70ef9db22d0e5p+2])}
is $\itv{0}{5.764625}$, a zero lower bound of the other sign than the vector code gives: with a $J$ from $1$, which the vector code takes, the result is $\itv{-0}{\code{0x1.e2ab041732196p-5}}$.

\medskip
The two paths of the SSE2 intervals, with and without the AVX-512 path
(\pxref{sec:avx512}), give the same bits, but
for the sign of a zero bound of \code{x * d}, \code{x *= d} and
\code{d * x} (computed as \code{x * d}) where the zero bounds of $x$ have
different signs, $\itv{+0}{-0}$: it is $\itv{+0}{-0}$ with SSE2 and
$\itv{-0}{+0}$ with AVX-512. The NEON intervals take the larger of four
products too,\newinvfive{} with the maximum of 64-bit ARM, which ranks
$+0$ above $-0$, where the \code{maxpd} of the SSE2 intervals returns its
second operand when the two compare equal: the lower bound of
\code{[-1, 0] * [-1, 0]} is $-0$ with the NEON intervals, where the SSE2
ones give $+0$ and the FPU ones the compiler's, and so may another zero
bound of a product. \code{interval::zero()} is stored
$\langle+0, +0\rangle$ as with SSE2, so that its lower bound is $-0$
(\pxref{sec:kernels}). The tool \file{tests/tools/arith} compares the NEON
and the FPU intervals on a grid of operands, with the calling code in each
of the four rounding directions: it found no other difference.

\section{Discussion}
%%-----------------

This chapter documents the tightness of each operation of \code{interval},
as 12.10.3 of IEEE 1788-2015 requires, from its algorithm, and gives the
largest distance the tests found. The basic operations, the integer
functions and the absmax functions are tightest, as the standard requires.
The elementary functions, the recommended functions and the error
functions are tightest too, where the standard requires validity and
recommends accuracy, because CORE-MATH is correctly rounded in the
direction in effect: the value at a bound and the double below it are the
tightest bounds, and the cases where the exact value is a double are told
apart with integers. The exceptions are stated with their size: the roots
of an order other than 1, 2 and 3 are within 2 doubles, their inverses
within 4, the integer power is one double beyond where $x^n$ is within $2K$ of a double, which is rare but not near $1$, and \code{compoundm1} in about one pair in $10^{17}$.
The reverse functions are valid, and accurate or within a few doubles of the
accurate bounds, as the chapter says of each, which the standard recommends,
and the tests measure how far they are from the values they keep.

\medskip
Several choices trade a little tightness or time for a proof. The real
power proves that its lower bound is a double with integers rather than
call CORE-MATH downward, which would double its cost; \code{compoundm1}
settles most pairs in a fast phase of 64-bit fixed point whose error bound
is proved, and sends the others to balls on 128 bits; the roots of an
order above 3 are proved by the power of the double on each side, at the
price of up to 2 doubles. The bounds are the same on every machine and in
every build, CORE-MATH being compiled into the library and the kernels
computing the integer powers with the same code; only the sign of a zero
bound differs between the SSE2, the NEON and the FPU intervals, and
between the compilers for the FPU ones, which changes no set.

\medskip
What remains is stated as such: the widths of \code{tan} between
$\underline{\pi}$ and $\pi$ are a theoretical case where the result is wider
than the tightest one, and the intervals of floats are not covered. The next chapter
describes how GAOL computes these bounds: the register that holds the
negated lower bound, the SSE2, AVX-512, NEON and FPU kernels, the check of
the rounding direction and its cost, and the comparison with the other
libraries.
